\documentclass[11pt]{article}
\usepackage[margin=2.2cm]{geometry}
\usepackage{booktabs}
\usepackage{array}
\usepackage[hidelinks]{hyperref}
\usepackage{xcolor}
\usepackage{listings}
\lstset{basicstyle=\ttfamily\small, breaklines=true, columns=fullflexible, frame=single, framerule=0.2pt}

\title{PATSTAT metric chain: cancellation record (2026-09-09)}
\author{Dawoon Jeong \\ \small \texttt{Science of Science/PATSTAT}}
\date{Written 2026-09-09 09:45 CDT}

\begin{document}
\maketitle

\section*{Summary}

The nine-job PATSTAT metric chain submitted on 2026-09-07 15:18 by
\texttt{jobs/PATSTAT/submit\_patstat.sh} was cancelled in full on
2026-09-09 09:41:08 CDT with \texttt{scancel}, at the user's request.
Only the first job, \texttt{ps\_patstat\_reference}, had started; it had run for
24\,h\,06\,m of its 48\,h limit without leaving cell~2 of the notebook.
The eight downstream jobs were still \texttt{PENDING (Dependency)} and never ran.
No output parquet was written, so there is nothing partial to clean up; the chain can be
resubmitted from scratch once the stall in \texttt{patstat\_reference} is understood.

\section*{1. Jobs cancelled}

All jobs ran on partition \texttt{jevans}, account \texttt{pi-jevans}, QOS \texttt{jevans},
16 CPUs and 400\,GB per node, 48\,h limit (\texttt{jobs/PATSTAT/nb.sbatch}).

\begin{center}
\small
\begin{tabular}{llllll}
\toprule
Job ID & Notebook & Depends on & State at cancel & Started & Elapsed \\
\midrule
57940901 & patstat\_reference            & (none)   & RUNNING & 2026-09-08 09:34:36 & 1-00:06:32 \\
57940903 & patstat\_metadata             & 57940901 & PENDING & --- & 0 \\
57940905 & patstat\_citation             & 57940903 & PENDING & --- & 0 \\
57940907 & patstat\_disruption           & 57940903 & PENDING & --- & 0 \\
57940911 & patstat\_sb                   & 57940903 & PENDING & --- & 0 \\
57940913 & patstat\_z\_score             & 57940903 & PENDING & --- & 0 \\
57940916 & patstat\_citation\_trend      & 57940905 & PENDING & --- & 0 \\
57940920 & patstat\_hit\_probability     & 57940905 & PENDING & --- & 0 \\
57940922 & patstat\_feg\_disruption\_trend & 57940907 & PENDING & --- & 0 \\
\bottomrule
\end{tabular}
\end{center}

All nine now show \texttt{CANCELLED by 30299883} in \texttt{sacct}. The stage-0 load array
(57939749, 21 tasks, 2026-09-07) is unaffected: \texttt{raw/} holds all 11 tables
(tls201, 202, 204, 206, 207, 209, 211, 212, 224, 228, 230; 22\,GB of zstd parquet) and remains valid input.

\section*{2. State of \texttt{ps\_patstat\_reference} (57940901) at cancellation}

\paragraph{Where it stopped.}
The log \texttt{jobs/PATSTAT/logs/ps\_patstat\_reference-57940901.out} ends with the
citation audit table that cell~2 prints \emph{before} it builds the edge table. The line
``patent citations routed to a citing application: \ldots'' that follows the
\texttt{CREATE OR REPLACE TABLE edges} statement never appeared, and cell~3 never began.
So the whole 24 hours were spent inside one DuckDB statement: the three-way join

\begin{lstlisting}
CREATE OR REPLACE TABLE edges AS
  SELECT ... FROM raw/tls212 c
  JOIN pub pc      ON c.pat_publn_id = pc.pat_publn_id
  LEFT JOIN pub pp ON c.cited_pat_publn_id <> 0
                  AND c.cited_pat_publn_id = pp.pat_publn_id
  WHERE (c.cited_pat_publn_id <> 0 OR c.cited_appln_id <> 0)
    AND NOT <third-party predicate>
\end{lstlisting}

(An earlier same-day summary attributed the stall to cell~3; the log shows it was cell~2.)

\paragraph{What the notebook had established before stalling.}
\begin{center}
\small
\begin{tabular}{lr}
\toprule
Quantity & Value \\
\midrule
Universe (\texttt{ipr\_type='PI'}, \texttt{appln\_id} $<$ 900\,000\,000, filing year 1900--2023) & 84\,914\,026 applications \\
Publications (\texttt{tls211}, \texttt{pat\_publn\_id} $\neq$ 0) & 151\,566\,588 \\
\texttt{tls212} rows total & 507\,520\,117 \\
\quad of which patent citations (cited publn or appln $\neq$ 0) & 434\,105\,665 \\
\quad of which NPL or empty & 73\,414\,452 \\
\quad of which third party (115 / TPO) & 132\,354 \\
\quad of which replenished & 17\,728\,407 \\
\bottomrule
\end{tabular}
\end{center}

Cells 0--2 up to the audit table completed within the first minute (log mtime 09:34).

\paragraph{Resource picture (sstat / ps on midway3-0290 just before cancel).}
\begin{center}
\small
\begin{tabular}{ll}
\toprule
Node & midway3-0290 (16 CPUs, 400\,GB requested) \\
Wall time & 24\,h\,06\,m of 48\,h \\
AveCPU & 16\,d\,00\,h\,33\,m $\approx$ 16 cores fully busy for the whole run \\
Python RSS at cancel & 39.6\,GB; MaxRSS 65.3\,GB \\
DuckDB memory\_limit & 200\,GB (\texttt{ps.connect}); never approached \\
MaxDiskRead / MaxDiskWrite & 5.0\,GB / 1.4\,kB \\
\texttt{cache/duckdb\_tmp} & empty --- no spill to disk \\
\texttt{output/} & empty --- \texttt{patstat\_reference.parquet} never written \\
\bottomrule
\end{tabular}
\end{center}

The statement was CPU-bound in memory, not I/O- or memory-bound: all threads were busy,
memory stayed at a third of the limit, nothing was spilled, and nothing was written.
A hash join of 434\,M rows against a 152\,M-row table on 16 threads should finish in minutes,
so a 24\,h run points to a pathological plan rather than to the data volume.
Software: DuckDB 1.5.5, pyarrow 22.0.0, Python 3.12.12 (env \texttt{Curvature}).

\section*{3. Hypotheses to test before resubmitting}

None of these were verified; the job was cancelled without an \texttt{EXPLAIN ANALYZE}.

\begin{enumerate}
  \item \textbf{The LEFT JOIN condition.} \texttt{c.cited\_pat\_publn\_id <> 0 AND c.cited\_pat\_publn\_id = pp.pat\_publn\_id}
        mixes a left-side filter into the join predicate. If the planner does not lift the
        \texttt{<> 0} out, the join may degrade from a hash join to a nested-loop or a
        filtered cross product. Rewrite with the filter in a CTE on \texttt{c} and a pure
        equality join, or use \texttt{NULLIF(c.cited\_pat\_publn\_id, 0) = pp.pat\_publn\_id}.
  \item \textbf{Two probes of the same 152\,M-row build side.} \texttt{pub} is joined twice.
        Check with \texttt{EXPLAIN ANALYZE} on one \texttt{tls212} part whether the plan builds
        the hash table once or once per probe, and whether the build side is \texttt{pub}
        (as intended) or \texttt{tls212}.
  \item \textbf{Parquet scan of \texttt{raw/tls212/*.parquet} inside the join.}
        Materialising \texttt{tls212} into a DuckDB table first (or the filtered projection
        only, about 434\,M rows and six columns) separates scan cost from join cost.
  \item \textbf{Timing on a subset.} Run cell~2 on a single part
        (\texttt{raw/tls212/part01.parquet}) on an interactive node with \texttt{\%\%time}
        and \texttt{EXPLAIN ANALYZE}; the full statement should scale roughly linearly from that.
\end{enumerate}

\section*{4. How to resubmit}

Nothing on disk needs deleting. After fixing cell~2 of
\texttt{notebook/patstat\_reference.ipynb} and confirming the subset timing:

\begin{lstlisting}
cd "/project/jevans/Dawoon/Science of Science/jobs/PATSTAT"
./submit_patstat.sh plan      # print the order
./submit_patstat.sh           # submit the nine-job chain (afterok throughout)
\end{lstlisting}

To reproduce the accounting in this note:

\begin{lstlisting}
sacct -j 57940901,57940903,57940905,57940907,57940911,57940913,57940916,57940920,57940922 \
      -X -P --format=JobID,JobName%32,State,Start,End,Elapsed
\end{lstlisting}

\end{document}
